% Generated draft: requires mathematical and source review.
\documentclass[11pt]{article}
\usepackage[margin=25mm]{geometry}
\usepackage{fontspec}
\setmainfont{Latin Modern Roman}
\usepackage{amsmath,amssymb,mathtools,unicode-math}
\setmathfont{Latin Modern Math}
\usepackage{enumitem,needspace,xurl,hyperref}
\setlength{\parindent}{0pt}
\setlength{\parskip}{5pt}
\setlength{\emergencystretch}{3em}
\newcommand{\N}{\mathbb{N}}
\newcommand{\Z}{\mathbb{Z}}
\newcommand{\Q}{\mathbb{Q}}
\newcommand{\R}{\mathbb{R}}
\newcommand{\C}{\mathbb{C}}
\newcommand{\F}{\mathbb{F}}
\newcommand{\A}{\mathbb{A}}
\newcommand{\PP}{\mathbb{P}}
\newcommand{\E}{\mathbb{E}}
\newcommand{\eps}{\varepsilon}
\newcommand{\dd}{\,\mathrm{d}}
\begin{document}
\section*{UnsolvedMath ID: 30001510}
\textbf{Source problem code:} catalogue entry
\subsection{Definitions and mathematical statement}
Let \(G\) be a free group, and write \([G,G]\) for its commutator subgroup, generated by the commutators \([a,b]=aba^{-1}b^{-1}\). For \(h\in[G,G]\), its commutator length \(\operatorname{cl}_G(h)\) is the least number of commutators whose product is \(h\), with \(\operatorname{cl}_G(1)=0\). Its stable commutator length is
\[
\operatorname{scl}_G(h)=\lim_{n\to\infty}\frac{\operatorname{cl}_G(h^n)}{n}.
\]
A quasimorphism is a function \(\phi:G\to\R\) with finite defect
\[
D(\phi)=\sup_{x,y\in G}|\phi(xy)-\phi(x)-\phi(y)|<\infty.
\]
It is homogeneous if \(\phi(x^n)=n\phi(x)\) for all \(x\in G\) and \(n\in\Z\). Under the usual positive-defect convention in stable commutator length, a homogeneous quasimorphism is extremal for \(g\in[G,G]\) if
\[
D(\phi)>0,\qquad \frac{\phi(g)}{2D(\phi)}=\operatorname{scl}_G(g).
\]
\textbf{Source question.} For every free group \(G\) and every \(g\in[G,G]\), does there exist an extremal quasimorphism for \(g\) arising from symplectic geometry?

To express the intended restriction, let \(\mathcal S_G\) stand for the class of quasimorphisms on \(G\) meant by ``arising from symplectic geometry.'' If homogeneous quasimorphisms with positive defect are intended, the question has the schematic form
\[
\forall\,G\text{ free}\quad\forall\,g\in[G,G]\quad
\exists\,\phi\in\mathcal S_G:\quad
\phi\text{ is homogeneous},\quad D(\phi)>0,\quad
\phi(g)=2D(\phi)\operatorname{scl}_G(g).
\]
The supplied record does not define \(\mathcal S_G\), specify the permitted symplectic constructions, or explicitly specify the convention for extremality. Thus this display records a standard interpretation, not a fully determined reformulation. The source gives no rank restriction and does not exclude the identity. For free groups of rank zero or one, every homogeneous quasimorphism has zero defect, so the positive-defect interpretation requires a qualification absent from the source. Allowing zero defect in the equality alone would make the zero quasimorphism satisfy it for every \(g\); that equality therefore cannot by itself resolve the intended boundary convention.
\subsection{Short English statement}
The question asks whether each element of the commutator subgroup of a free group admits a quasimorphism that realizes its stable commutator length and comes from symplectic geometry. The permitted geometric constructions and the treatment of zero defect and low-rank free groups are unspecified. The supplied literature assessment also discusses every extremal quasimorphism, a stronger quantifier than the existence question stated here.
\subsection{Sources}
\begin{itemize}
\item Local UnsolvedMath registry, supplied complete record: permanent ID \(30001510\), exact source problem code \texttt{catalogue entry}, title ``Symplectic Extremal Quasimorphisms on Free Groups,'' in the Oberwolfach Reports Open Problems collection. The proposed year is recorded as 2010; no individual proposer is supplied.
\item Supplied source citation: \emph{Geometric Group Theory, Hyperbolic Dynamics and Symplectic Geometry} (2011), Oberwolfach Reports, pp. 1761--1825, DOI \texttt{10.4171/owr/2010/30}. \url{https://doi.org/10.4171/owr/2010/30}. The background describes the report as documenting workshop discussions and open problems. This authoritative report is cited by the record; it has not been independently consulted for this draft.
\item Supplied literature source labelled ``Calegari Walker'': \url{https://arxiv.org/abs/1212.2618}. The record supplies no title or precise result identifying how this reference bears on the question.
\item Supplied literature source labelled ``original OWR report'': \url{https://doi.org/10.4171/owr/2010/30}.
\item Literature assessment dated exactly \texttt{2026-08-21}: the supplied status is \texttt{partially\_solved}. The assessment states that free groups have explicit extremal quasimorphisms and broad algorithmic constructions, while the stronger requirement that every extremal quasimorphism arise from symplectic geometry is established only in special families. This is supplied dated curation derived from the OWR record, not independently verified here; its time-sensitive status and cited claims require independent verification. Its universal claim about extremal quasimorphisms differs from the existence question in the statement.
\item The supplied statement status is \texttt{corrected\_verified}. Its accompanying verification note reports comparison with the official Oberwolfach report and normalization of notation, context, and question form without changing the source-level claim. These are supplied verification assertions, not independent verification performed for this draft.
\end{itemize}
\end{document}
